%% ===========================================================================
%%  2. Related Work
%% ===========================================================================

\section{Related Work}
\label{sec:related}

Our work sits at the intersection of three research areas: (i)~trust and
explainability in human-AI interaction, (ii)~self-supervised audio
representation learning for music, and (iii)~conformal prediction for
distribution-free uncertainty quantification. We review each in turn and
identify the gap that our framework fills.

\subsection{Trust and Explainability in Human-AI Interaction}

A central challenge in human-AI interaction is that users cannot assess the
reliability of AI outputs when only point predictions are
provided~\cite{amershi2019guidelines}. This trust gap is particularly acute in
creative and perceptual domains, where ground truth is ambiguous and system
errors may go unnoticed by non-experts. Miller~\cite{miller2019explanation}
argues that explanations in AI should be contrastive, selective, and social---in
other words, they should help users understand \emph{why} a particular output
was produced rather than \emph{how} the model works internally.

Existing approaches to communicating model confidence fall into two broad
categories. \emph{Post-hoc explanation} methods, such as
LIME~\cite{ribeiro2016lime}, highlight input features that influenced a
prediction but provide no statistical guarantee about the prediction's
correctness. \emph{Bayesian uncertainty} techniques, including Monte Carlo
dropout~\cite{gal2016dropout} and deep ensembles, produce variance estimates
that are difficult for non-technical users to interpret and offer only
asymptotic calibration.

Our work proposes a third path: instead of explaining \emph{why} a model chose a
label, we communicate \emph{how many} labels could plausibly be correct. The
prediction set size---a single integer---provides an immediately interpretable
confidence signal that requires no understanding of probability, entropy, or
model architecture. This aligns with the principle articulated by Amershi et
al.~\cite{amershi2019guidelines} that AI systems should ``convey the
consequences of user actions'' and ``support efficient dismissal'' of
suggestions.

\subsection{Music AI and Creative Workflows}

The adoption of AI in music tools has grown rapidly, from style
transfer~\cite{huang2020creative} to generative composition~\cite{copet2023musicgen}
to automated playlist curation. In these applications, genre classification
serves as a fundamental building block for downstream tasks such as filtering,
tagging, and recommendation. However, music genre is inherently subjective and
multidimensional---a single track may legitimately belong to multiple genres,
and annotators frequently disagree~\cite{tzanetakis2002gtzan}. Traditional
single-label classifiers force this ambiguity into a single output, misleading
users about the certainty of the classification.

For generative music specifically, evaluation remains an open
challenge~\cite{yang2018evaluation}. Metrics such as the Fr\'{e}chet Audio
Distance (FAD)~\cite{kilgour2019fad} measure distribution-level similarity but
provide no per-sample feedback. A producer cannot ask ``is \emph{this
particular} generated track convincing as jazz?'' and receive an answer with
quantified confidence. Our prediction sets address this need by providing
instance-level reliability information.

\subsection{Self-Supervised Audio Representations}

Self-supervised learning (SSL) has become the dominant paradigm for extracting
features from raw audio. In the speech domain, Wav2Vec~2.0~\cite{baevski2020wav2vec2}
and HuBERT~\cite{hsu2021hubert} demonstrated that learned representations
surpass hand-crafted features on a wide range of tasks. For music,
MERT~\cite{li2023mert} extends this approach: a 24-layer Transformer
pre-trained on 160,000 hours of music via masked language modeling with dual
teacher supervision (acoustic RVQ-VAE and musical CQT). MERT produces
1024-dimensional embeddings that encode both timbral and harmonic information,
and has been shown to support strong performance across MIR benchmarks when
paired with task-specific models.

In this work, we use MERT-v1-330M purely as a frozen feature extractor,
extracting embeddings from the last hidden layer via mean-pooling over time. We
then train a lightweight Linear Probe on these embeddings, keeping the
underlying representation fixed. This design choice keeps training fast and
reproducible while leveraging the rich acoustic knowledge embedded in MERT's
pre-trained weights. For a detailed comparison of audio SSL models and
alternative multimodal approaches such as CLAP, we refer the reader to
comprehensive surveys in the MIR literature.

\subsection{Conformal Prediction as a User-Facing Interface}

Conformal prediction (CP)~\cite{vovk2005algorithmic,shafer2008tutorial} is a
distribution-free, finite-sample framework for constructing prediction sets with
formal coverage guarantees: given a significance level $\alpha$, CP produces a
set $\Gamma(x)$ that contains the true label with probability at least
$1-\alpha$, under only the assumption that calibration and test data are
exchangeable. A modern introduction is provided by Angelopoulos and
Bates~\cite{angelopoulos2023gentle}.

CP has primarily been studied as a statistical methodology for machine learning
practitioners. Applications include regression~\cite{romano2019conformalized},
image classification with class-conditional
sets~\cite{angelopoulos2021uncertainty}, and software libraries for practical
deployment~\cite{cordier2023mapie}. However, to our knowledge, no prior work has
explicitly positioned CP as a \emph{user interface layer}---a mechanism through
which AI systems communicate their certainty to non-expert users in an
interpretable, actionable format. This is the perspective we adopt.

From an HCI standpoint, CP has several attractive properties as an interface
primitive:
\begin{enumerate}[leftmargin=*,itemsep=0pt]
    \item \textbf{Set size as a natural confidence indicator.} Unlike
    probabilities or entropy values, the cardinality of a prediction set is
    trivially interpretable: smaller means more confident, larger means more
    uncertain. This maps to the long-established finding that humans reason more
    reliably about categorical sets than about continuous probability
    values~\cite{miller2019explanation}.
    \item \textbf{$\alpha$ as a direct user control.} The significance level
    functions as a slider that users can adjust according to their risk
    tolerance, without needing to understand the underlying calibration
    mechanism.
    \item \textbf{Distribution-free guarantees.} The coverage guarantee holds
    regardless of model architecture, embedding space geometry, or score
    distribution shape, making the interface behavior predictable and uniform
    across different underlying systems.
\end{enumerate}

Our contribution is to realize this interface perspective concretely for music
genre classification, demonstrating that CP not only provides theoretical
guarantees but also yields practically useful prediction sets that can inform
user decision-making in creative workflows.
